% TOP_PROBLEM: {"id": 20001094, "title": "Quantitative Hales-Jewett bounds and exact binary calibration", "category_id": 2, "category": {"id": 2, "name": "combinatorics", "display_name": "Combinatorics", "description": "Counting problems, graph theory, discrete structures.", "slug": "combinatorics", "order_index": 2, "created_at": "2026-07-31T15:26:25.671Z"}, "status": "partially_solved", "problem_number": "AIM-COMBINATORICS-0219", "set_id": 14, "statusQualification": "The source specifies no numerical conjecture; both minimal thresholds and their parameter regimes are defined explicitly."}
% FIRST_POSTED: 2026-10-01
% ENTRY_KIND: statement_only
% UnsolvedMath ID 20001094; source catalogue code AIM-COMBINATORICS-0219
% !TeX program = lualatex
% Definition and sources only; this file is not an LLM solution attempt.
\documentclass[11pt,a4paper]{article}
\usepackage[margin=25mm]{geometry}
\usepackage{fontspec}
\setmainfont{lmroman10-regular.otf}[BoldFont=lmroman10-bold.otf,
 ItalicFont=lmroman10-italic.otf,BoldItalicFont=lmroman10-bolditalic.otf]
\usepackage{amsmath,amssymb,mathtools}
\usepackage{unicode-math}
\setmathfont{latinmodern-math.otf}
\AtBeginDocument{\let\setminus\smallsetminus}
\usepackage{xurl,hyperref}
\hypersetup{colorlinks=true,urlcolor=blue}
\setcounter{secnumdepth}{0}
\setlength{\parindent}{0pt}
\setlength{\parskip}{5pt}
\setlength{\emergencystretch}{3em}
\begin{document}
\section*{Quantitative Hales-Jewett bounds and exact binary calibration}\label{20001094}
{\small UnsolvedMath ID 20001094 · AIM-COMBINATORICS-0219 · Combinatorics · AIM Workshop Problem Lists}
\subsection{Definitions and mathematical statement}
For integers $q\ge2$ and $n\ge1$, let $[q]^n$ be the words of length
$n$ over the alphabet $[q]=\{1,\ldots,q\}$. A combinatorial line
is obtained by choosing a nonempty set of coordinate positions $I$,
fixing one letter at every position outside $I$, and making all
coordinates in $I$ equal to the same variable letter. Letting that
letter run through $[q]$ gives the $q$ words of the line.
The variable positions need not be consecutive.

Define $\operatorname{HJ}(q,r)$, for $r\ge2$, as the least $n$ such
that every coloring $[q]^n\to[r]$ contains a monochromatic
combinatorial line. For $0<\delta\le1$, define
$\operatorname{DHJ}(q,\delta)$ as the least $n$ such that every
$A\subseteq[q]^n$ satisfying $|A|\ge\delta q^n$ contains a
combinatorial line. These are respectively the coloring and density
Hales--Jewett thresholds.

The workshop asks for substantially improved, explicit quantitative
bounds for both thresholds, ideally approaching their correct growth
in the alphabet size and coloring or density parameter. In particular,
fixed $q$ with $r\to\infty$, and fixed $q$ with $\delta\downarrow0$,
are meaningful asymptotic regimes. Mere finiteness does not answer
the quantitative question. The source's word ``reasonable'' specifies
no particular numerical conjecture, so no arbitrary growth bound is
inserted here. A density line must contain all $q$ words, not merely
two words agreeing on some coordinates.
\subsection{Short English statement}
Obtain sharp usable quantitative bounds for both the coloring and density forms of the Hales–Jewett theorem.
\subsection{Sources}
\begin{itemize}
\item[S1] ULAM.AI and the UnsolvedMath Contributors, supplied problems.json, record 20001094 (AIM-COMBINATORICS-0219), AIM Workshop Problem Lists. Catalogue identity and source transcription; CC BY 4.0.
\url{https://huggingface.co/datasets/ulamai/UnsolvedMath}.
\item[S2] American Institute of Mathematics, Recent trends in additive combinatorics, item 1.13. Original workshop question underlying AIM-COMBINATORICS-0219.
\url{https://aimath.org/WWN/additivecomb/additivecomb.pdf}.
\end{itemize}
\end{document}
